% Template for post or presentation abstracts submitted to ILAS 2022
% 1. Fill in the presenter's information (name, affiliation, email
% address) and talk information (title of presentation, abstract).
% 2. Compile to make sure there are no errors.
% 3. Save the .tex file for your abstract with a distinctive name,
% ideally "FamilynameGivenname.tex".
% 4. Upload your .tex file to https://clr.ie/129557
%       
% * Please keep your LaTeX commands simple, and avoid the use of
%    user-defined macros.
% * Include details of co-authors and funding acknowledgements after the abstract.
% * Upload only the LaTeX file (with .tex extension).
% * Questions: Contact Niall Madden (Niall.Madden@NUIGalway.ie)

\documentclass[a4paper]{amsart} 
\usepackage{amssymb} 
\usepackage{hyperref}
\pagestyle{empty}
\date{} 

%% IGNORE THE NEXT 4 LINES
\makeatletter % 
\def\labelenumi{\theenumi.}
\def\theenumi{\@arabic\c@enumi}
\makeatother
%% STOP IGNORING NOW

\begin{document}

\title{On a conjecture by Graham and Lov\'{a}sz concerning the distance characteristic polynomial of block graphs}
\author{Aida Abiad} %Presenting author only
\address{Eindhoven University of Technology, Ghent University, Vrije universiteit Brussel} % Affiliation of presenting author
\email{a.abiad.monge@tue.nl} % Email address of presenting author only

\maketitle

% The abstract  ***no more than one page***

Graham and Lov\'{a}sz conjectured in 1978 that the sequence of normalized coefficients of the distance characteristic polynomial of a tree is unimodal with the maximum value occurring at $\lfloor\frac{n}{2}\rfloor$ for a tree $T$ of order $n$. In this paper we extend this old conjecture to block graphs. In particular, we prove the unimodality part and we establish the peak for several extremal cases of block graphs.  \\
This is joint work with B. Brimkov, A. Khramova, J. Koolen and S. Hayat.


\end{document}


% Please leave the next bit alone  
%%% Local Variables: 
%%% mode: latex
%%% TeX-master: "../ILAS-2022"
%%% End:
